\documentclass[10pt,a4paper]{amsart}
\usepackage{amsmath,amssymb,amsthm,graphicx}
% Retained minimal preamble for the Stacks Project chapter more-algebra.tex (More on Algebra).
% The chapter file pinned-original-source.tex carries only \input{preamble} on its first line, so no
% preamble text lives inside the chapter itself.  The Stacks Project uses one single shared preamble for
% every chapter; the snapshot stages it identically as sources/stacks-chow/preamble.tex and
% sources/stacks-etale-cohomology/preamble.tex.  The block below is copied mechanically, line for line,
% from sources/stacks-chow/preamble.tex lines 14-18 (the Xy-pic setup) and 154-217 (the native theorem
% declarations together with the mathematical macros).  Nothing was added, removed or redefined; the
% wrapper class and the amsmath/amssymb/amsthm/graphicx base packages are supplied by the assembler.
% Chapter: more-algebra.tex -- More on Algebra.
% For commutative diagrams we use Xy-pic
\usepackage[all]{xy}

% We use 2cell for 2-commutative diagrams.
\xyoption{2cell}

% Theorem environments.
%
\theoremstyle{plain}
\newtheorem{theorem}[subsection]{Theorem}
\newtheorem{proposition}[subsection]{Proposition}
\newtheorem{lemma}[subsection]{Lemma}

\theoremstyle{definition}
\newtheorem{definition}[subsection]{Definition}
\newtheorem{example}[subsection]{Example}
\newtheorem{exercise}[subsection]{Exercise}
\newtheorem{situation}[subsection]{Situation}

\theoremstyle{remark}
\newtheorem{remark}[subsection]{Remark}
\newtheorem{remarks}[subsection]{Remarks}

\numberwithin{equation}{subsection}

% Macros
%
\def\lim{\mathop{\mathrm{lim}}\nolimits}
\def\colim{\mathop{\mathrm{colim}}\nolimits}
\def\Spec{\mathop{\mathrm{Spec}}}
\def\Hom{\mathop{\mathrm{Hom}}\nolimits}
\def\Ext{\mathop{\mathrm{Ext}}\nolimits}
\def\SheafHom{\mathop{\mathcal{H}\!\mathit{om}}\nolimits}
\def\SheafExt{\mathop{\mathcal{E}\!\mathit{xt}}\nolimits}
\def\Sch{\mathit{Sch}}
\def\Mor{\mathop{\mathrm{Mor}}\nolimits}
\def\Ob{\mathop{\mathrm{Ob}}\nolimits}
\def\Sh{\mathop{\mathit{Sh}}\nolimits}
\def\NL{\mathop{N\!L}\nolimits}
\def\CH{\mathop{\mathrm{CH}}\nolimits}
\def\proetale{{pro\text{-}\acute{e}tale}}
\def\etale{{\acute{e}tale}}
\def\QCoh{\mathit{QCoh}}
\def\Ker{\mathop{\mathrm{Ker}}}
\def\Im{\mathop{\mathrm{Im}}}
\def\Coker{\mathop{\mathrm{Coker}}}
\def\Coim{\mathop{\mathrm{Coim}}}

% Boxtimes
%
\DeclareMathSymbol{\boxtimes}{\mathbin}{AMSa}{"02}

%
% Macros for moduli stacks/spaces
%
\def\QCohstack{\mathcal{QC}\!\mathit{oh}}
\def\Cohstack{\mathcal{C}\!\mathit{oh}}
\def\Spacesstack{\mathcal{S}\!\mathit{paces}}
\def\Quotfunctor{\mathrm{Quot}}
\def\Hilbfunctor{\mathrm{Hilb}}
\def\Curvesstack{\mathcal{C}\!\mathit{urves}}
\def\Polarizedstack{\mathcal{P}\!\mathit{olarized}}
\def\Complexesstack{\mathcal{C}\!\mathit{omplexes}}
% \Pic is the operator that assigns to X its picard group, usage \Pic(X)
% \Picardstack_{X/B} denotes the Picard stack of X over B
% \Picardfunctor_{X/B} denotes the Picard functor of X over B
\def\Pic{\mathop{\mathrm{Pic}}\nolimits}
\def\Picardstack{\mathcal{P}\!\mathit{ic}}
\def\Picardfunctor{\mathrm{Pic}}
\def\Deformationcategory{\mathcal{D}\!\mathit{ef}}

% --- appended 2026-10-12 (Chief): the project own definitions of the reference/slogan/history
% display environments, copied mechanically from the shared Stacks preamble lines 8-12
% (\usepackage{verbatim} plus the three \newenvironment lines, which the chapter preamble keeps
% verbatim from the project preamble). Nothing added, removed or redefined: these environments
% carry no mathematical content and the project itself typesets them as comments. Without these
% three lines a counted statement whose own text uses them fails with Environment slogan undefined.
\usepackage{verbatim}
\newenvironment{reference}{\comment}{\endcomment}
\newenvironment{slogan}{\comment}{\endcomment}
\newenvironment{history}{\comment}{\endcomment}
\usepackage{cleveref}
\crefname{originallemma}{original lemma}{original lemmas}
\Crefname{originallemma}{Original lemma}{Original lemmas}
\makeatletter
\newlabel{lemma-map-out-of-almost-free}{{22338}{source}{}{sourcepointer-1}{}}
\newlabel{lemma-map-out-of-almost-free@cref}{{[originallemma][0][]22338}{[1][0][]source}{}{}{}}
\makeatother
\crefname{originallemma}{original lemma}{original lemmas}
\Crefname{originallemma}{Original lemma}{Original lemmas}
\makeatletter
\newlabel{lemma-represent-two-term-complex}{{22304}{source}{}{sourcepointer-2}{}}
\newlabel{lemma-represent-two-term-complex@cref}{{[originallemma][0][]22304}{[1][0][]source}{}{}{}}
\makeatother
\title{Annihilation of Ext one by an ideal is equivalent to the multiplication maps on the first term factoring through the differential}
\author{The Stacks Project Authors}
\date{stacks-project:more-algebra.tex:local-snapshot-20261012; GNU FDL 1.2 or later}
\begin{document}
\maketitle
\setcounter{subsection}{1}
\noindent\textit{Editorial note.} All delivered mathematical text is the author's own, copied from the pinned Stacks Project chapter file more-algebra.tex; no statement or proof step was rewritten, reordered or removed.\par\medskip
\noindent\textit{Editorial note.} The chapter file pinned-original-source.tex carries only \textbackslash{}input\{preamble\} on its first line, so no preamble text lives inside the chapter. The retained file sources/stacks-more-algebra/retained-preamble.tex is a mechanical, line-for-line copy of the project's single shared preamble as staged in the snapshot; nothing was added, removed or redefined and no package was added.\par\medskip
\noindent\textit{Editorial note.} The source URL is the source-level Stacks Project URL, as required when more than one card is produced from one source; the per-claim tag is carried by claim\_key.\par\medskip
\noindent\textit{Editorial note.} The counted proof span begins with a blank line, which is reproduced. The author writes the incomplete sentence 'Since is annihilated by a in I we see that we get the dotted arrow'; it is reproduced unchanged rather than silently repaired. The four references inside the counted proof span are recorded as non-load-bearing navigation pointers.\par\medskip
\noindent\textit{Allowed standard prerequisites (editorial).}\begin{itemize}
\item The description of maps out of a two-term complex in the derived category. invoked at proof line 22483. Lets an element of Ext\textasciicircum{}1\_R(K, N) be represented by a map from the first term, which the computation a varphi = varphi o h o d\_K\textasciicircum{}\{-1\} relies on.
\item The representation lemma for a complex concentrated in two degrees, with a projective or finite free first term. invoked at proof lines 22495 and 22506. Supplies the normal form of the representing complex in conditions (2) and (5).
\item Ext\textasciicircum{}1\_R(K, N) as classes of maps and the identification of the class of the identity on K\textasciicircum{}\{-1\}. used at proof lines 22484-22486. Is the choice that turns the annihilation hypothesis into the existence of the dotted arrow.
\end{itemize}\medskip
\ifdefined\hypertarget\hypertarget{sourcepointer-1}{}\fi
\noindent\textit{Source reference.} Original lemma 22338: sources/stacks-more-algebra/pinned-original-source.tex, source lines 22338--22338. Same-chapter navigation pointer; the statement it names is standard for the section and the argument that consumes it stays inside the counted proof span.\par\medskip
\ifdefined\hypertarget\hypertarget{sourcepointer-2}{}\fi
\noindent\textit{Source reference.} Original lemma 22304: sources/stacks-more-algebra/pinned-original-source.tex, source lines 22304--22304. Same-chapter navigation pointer; the statement it names is standard for the section and the argument that consumes it stays inside the counted proof span.\par\medskip
\begin{lemma}
\label{lemma-ext-1-annihilated-definite}
Let $R$ be a ring and let $I \subset R$ be an ideal.
Let $K \in D(R)$. Assume $H^i(K) = 0$ for $i \not \in \{-1, 0\}$.
The following are equivalent
\begin{enumerate}
\item $\Ext^1_R(K, N)$ is annihilated by $I$ for all $R$-modules $N$,
\item $K$ can be represented by a complex $K^{-1} \to K^0$
with $K^0$ free such that for any $a \in I$ the map
$a : K^{-1} \to K^{-1}$ factors through $d_K^{-1} : K^{-1} \to K^0$,
\item whenever $K$ is represented by a two term complex
$K^{-1} \to K^0$ with $K^0$ projective, then for any $a \in I$ the map
$a : K^{-1} \to K^{-1}$ factors through $d_K^{-1} : K^{-1} \to K^0$.
\end{enumerate}
If $R$ is Noetherian and $H^i(K)$ is a finite $R$-module for $i = -1, 0$,
then these are also equivalent to
\begin{enumerate}
\item[(4)] $\Ext^1_R(K, N)$ is annihilated by $I$ for every finite
$R$-module $N$,
\item[(5)] $K$ can be represented by a complex $K^{-1} \to K^0$
with $K^0$ finite free and $K^{-1}$ finite such that for any $a \in I$ the map
$a : K^{-1} \to K^{-1}$ factors through $d_K^{-1} : K^{-1} \to K^0$.
\end{enumerate}
\end{lemma}

\begin{proof}

\medskip\noindent
Assume (1) and let $K^{-1} \to K^0$
be a two term complex representing $K$
with $K^0$ projective.
We will use the description of maps in $D(R)$ out of $K^\bullet$ given in
Lemma \ref{lemma-map-out-of-almost-free} without further mention.
Choosing $N = K^{-1}$ consider the element $\xi$ of $\Ext^1_R(K, N)$
given by $\text{id}_{K^{-1}} : K^{-1} \to K^{-1}$. Since is annihilated
by $a \in I$ we see that we get the dotted arrow fitting into the
following commutative diagram
$$
\xymatrix{
K^{-1} \ar[d]_a \ar[r]_{d_K^{-1}} & K^0 \ar@{..>}[ld]^h \\
K^{-1}
}
$$
This proves that (3) holds. Part (3) implies (2) in view of
Lemma \ref{lemma-represent-two-term-complex} part (1). Assume
$K^\bullet$ is as in (2) and $N$ is an arbitrary $R$-module.
Any element $\xi$ of $\Ext^1_R(K, N)$ is given as the class of
a map $\varphi : K^{-1} \to N$. Then for $a \in I$ by assumption
we may choose a map $h$ as in the diagram above and we see that
$a\varphi = \varphi \circ a = \varphi \circ h \circ \text{d}_K^{-1}$
which proves that $a \xi$ is zero in $\Ext^1_R(K, N)$.
Thus (1), (2), and (3) are equivalent.

\medskip\noindent
Assume $R$ is Noetherian and $H^i(K)$ is a finite $R$-module for $i = -1, 0$.
Part (3) implies (5) in view of
Lemma \ref{lemma-represent-two-term-complex} part (2).
It is clear that (5) implies (2). Trivially (1) implies (4).
Thus to finish the proof it suffices to show that (4) implies
any of the other conditions. Let $K^{-1} \to K^0$ be a complex
representing $K$ with $K^0$ finite free and $K^{-1}$ finite as in
Lemma \ref{lemma-represent-two-term-complex} part (2). The argument
given in the proof of (2) $\Rightarrow$ (1)
shows that if $\Ext^1_R(K, K^{-1})$ is annihilated by $I$,
then (1) holds. In this way we see that (4) implies (1) and the
proof is complete.
\end{proof}

\end{document}
